% ============================================================
% CAS/Python ενότητες — Κεφάλαιο 9: Εφαρμογές Παραγώγου
%
% QR code: qr_1c92f8574017.png  (→ latex/qrcodes/)
% URL: https://nikmatz.github.io/ELADIC_GR/code/ch09/
% ============================================================

\casactivbox%
  {Maxima: Ακρότατα, Βελτιστοποίηση, Newton-Raphson}%
  {%
    \label{cas:ch9_deriv_apps}
    Δίνεται $f(x)=x^3-3x^2-9x+5$.
    \begin{enumerate}[label=(\alph*),itemsep=2pt]
      \item Βρείτε κρίσιμα σημεία ($f'=0$) με \texttt{solve}.
        Εφαρμόστε κριτήριο 2ης παραγώγου και χαρακτηρίστε τα
        ως τοπικό μέγιστο ή ελάχιστο.
      \item Βρείτε σημεία καμπής ($f''=0$) και προσδιορίστε
        διαστήματα κυρτότητας/κοιλότητας.
      \item \textbf{Βελτιστοποίηση}: ορθογώνιο με περίμετρο $P=20$ —
        βρείτε $x^*$ που μεγιστοποιεί το εμβαδόν $E(x)=x(10-x)$.
      \item \textbf{Θεώρημα Μέσης Τιμής}: για $f(x)=x^2-x+1$
        στο $[1,4]$, βρείτε $c$ τέτοιο ώστε
        $f'(c)=\dfrac{f(4)-f(1)}{4-1}$.
      \item \textbf{Newton-Raphson}: εκτελέστε 5 επαναλήψεις
        για εύρεση ρίζας $x^3-x-2=0$, ξεκινώντας από $x_0=1.5$.
    \end{enumerate}
    \smallskip
    \textbf{Εντολές Maxima:}
    \texttt{solve(diff(f,x)=0,x)}, \texttt{subst(c,x,diff(f,x,2))},
    \texttt{assume(x>0)}, \texttt{float()}
  }%
  {https://nikmatz.github.io/ELADIC_GR/code/ch09/}%
  {qr_1c92f8574017.png}

\subsection*{Δραστηριότητα Python}

\begin{pybox}{Python: Ακρότατα, Βελτιστοποίηση, Newton}%
             {https://nikmatz.github.io/ELADIC_GR/code/ch09/}%
             {qr_1c92f8574017.png}
\label{py:ch9_deriv_apps}
\begin{enumerate}[label=(\alph*),itemsep=2pt]
  \item Με \texttt{sympy.solve(diff(f,x),x)} βρείτε κρίσιμα
    σημεία του $f(x)=x^3-3x^2-9x+5$ και εφαρμόστε κριτήριο
    2ης παραγώγου. Σχεδιάστε $f$, $f'$, $f''$ σε κοινό διάγραμμα.
  \item Βελτιστοποίηση: επιλύστε το πρόβλημα ορθογωνίου
    με \texttt{sympy.solve} και επαληθεύστε αριθμητικά με
    \texttt{scipy.optimize.minimize\_scalar}.
  \item Για το κλειστό δοχείο ($V=32$): ελαχιστοποιήστε την
    επιφάνεια $S(x)=x^2+128/x$. Σχεδιάστε $S(x)$.
  \item Υλοποιήστε Newton-Raphson βήμα-βήμα για $x^3-x-2=0$
    και επαληθεύστε με \texttt{scipy.optimize.newton}.
    Σχεδιάστε τα βήματα σύγκλισης γραφικά.
\end{enumerate}
\medskip
\textbf{Βιβλιοθήκες / Εντολές:}
\texttt{sympy.solve()}, \texttt{sympy.lambdify()},
\texttt{scipy.optimize.minimize\_scalar()},
\texttt{scipy.optimize.newton()}, \texttt{matplotlib.pyplot}
\end{pybox}
